\chapter[AUTRES PAQUETS]{Autres paquets}\label{autres-paquets}

\section{Paquets invalides ou de
test}\label{paquets-invalides-ou-de-test}

DTI \texttt{,}. Ex : un Mic-E génère ça s'il détecte que la sentence GPS
reçue n'est pas valide.

\begin{verbatim}
, | Invalid Data or Test Data
\end{verbatim}

Ex :
\texttt{,191146,V,4214.2466,N,07303.5181,W,417.238,114.5,091099,14.7,W/GPS\ FIX}
--- GPS invalide d'un Mic-E qui a stripé le header \texttt{\$GPRMC} car
le \texttt{V} indique donnée invalide.

\section{Tous les autres paquets}\label{tous-les-autres-paquets}

Paquets ne correspondant à aucun format décrit = beacons non-APRS. Les
programmes peuvent les traiter ou les ignorer, mais doivent les traiter
sans effet secondaire. Souvent traités comme Status Reports (permet de
capter comme statut toute trame UI adressée à un beacon typique, ex :
trames d'ID de TNC). Une fois qu'un vrai Status Report (DTI
\texttt{\textgreater{}}) est reçu d'une station, il n'est plus écrasé
par ces autres paquets.

\begin{center}\rule{0.5\linewidth}{0.5pt}\end{center}
